\section{Appendix: Cost Model Details and Normalization}
\label{sec:cost_model}

This appendix details the acquisition cost model used by \textsc{VOILA}.
We describe the fidelity set, the empirical bandwidth proxy, the tier-aware
cost formulation, the normalization procedure, and the empirical measurement
methodology used to derive reported costs.

\subsection{Fidelity Set}
\label{sec:fidelity_set}

We consider a discrete set of input fidelities
\[
\mathcal{F}
=\{
\texttt{caption},
\texttt{resize}_{32\times32},
\texttt{jpeg\_q1},
\texttt{jpeg\_q10},
\texttt{full}
\},
\]
chosen to reflect representations that are persistently available in real-world
multimodal systems rather than artificially synthesized intermediates.
These fidelities correspond to natural storage and transmission breakpoints
across edge--cloud pipelines, agentic memory systems, and cyber--physical
deployments, where small artifacts are cached in low-latency tiers and
high-resolution images are retrieved from progressively colder storage tiers
with higher access cost and latency
\cite{aws_s3_storage_classes, gcs_storage_classes, azure_blob_storage}.

\subsection{Empirical Bandwidth Proxy}
\label{sec:bandwidth}

For each dataset, all fidelities in $\mathcal{F}$ are explicitly materialized
and their average on-disk size $\texttt{size}(f)$ is measured in kilobytes.
We define a dimensionless bandwidth ratio
\begin{equation}
r(f)
\;=\;
\frac{\texttt{size}(f)}{\texttt{size}(\texttt{full})},
\label{eq:ratio_def}
\end{equation}
which grounds bandwidth cost directly in empirical measurements and avoids
reliance on model-internal token counts or hardware-specific throughput
assumptions.

\subsection{Tier-Aware Acquisition Cost}
\label{sec:tier_cost}

Acquisition cost in practice is dominated by which storage or transmission tier
is accessed, with payload size acting as a secondary modifier. This abstraction
mirrors the pricing and performance models of major object storage systems,
where access tier determines latency and per-request cost, while object size
determines transfer charges
\cite{aws_s3_pricing, gcs_pricing, azure_blob_pricing}.
We define an unnormalized tier-aware cost
\begin{equation}
\tilde{c}(f)
\;=\;
b_{\mathrm{tier}}(f)
\;+\;
w_{\mathrm{bw}} \cdot r(f),
\label{eq:tier_cost}
\end{equation}
where $b_{\mathrm{tier}}(f)$ is a base cost determined by the access tier
(edge, warm, or cold), and $w_{\mathrm{bw}}$ controls the influence of
measured size differences within a tier.

\subsection{Normalization}
\label{sec:normalization}

For consistency across datasets and experiments, costs are normalized so that
full-resolution access has fixed cost 100:
\begin{equation}
c(f)
\;=\;
100 \cdot
\frac{\tilde{c}(f)}{\tilde{c}(\texttt{full})}.
\label{eq:normalize_100}
\end{equation}
All reported costs in the paper use this normalized value $c(f)$.
Only relative cost differences matter for VOI-based decision-making.

\subsection{Deployment Profiles}
\label{sec:deployment}

Table~\ref{tab:cost_weights} lists the tier base-costs and bandwidth weights used
for the three deployment profiles studied in this work.

\begin{table}[t]
\centering
\small
\caption{\textbf{Tier base-costs and bandwidth weights.}
Costs are computed via Eq.~\eqref{eq:tier_cost} and normalized using
Eq.~\eqref{eq:normalize_100}.}
\label{tab:cost_weights}
\begin{tabular}{lccccc}
\toprule
\textbf{Setting} &
$b_{\mathrm{tier}}(\texttt{caption})$ &
$b_{\mathrm{tier}}(\texttt{resize})$ &
$b_{\mathrm{tier}}(\texttt{jpeg\_q1})$ &
$b_{\mathrm{tier}}(\texttt{jpeg\_q10})$ &
$w_{\mathrm{bw}}$ \\
\midrule
Edge--Cloud VQA & 0.08 & 0.16 & 0.40 & 0.56 & 0.06 \\
Agentic Memory & 0.06 & 0.12 & 0.36 & 0.52 & 0.06 \\
CPS / IoT & 0.04 & 0.10 & 0.30 & 0.46 & 0.12 \\
\bottomrule
\end{tabular}
\end{table}

In all profiles, $\texttt{full}$ is treated as cold-tier access with
$b_{\mathrm{tier}}(\texttt{full}) = 1.00$, yielding $c(\texttt{full}) = 100$
after normalization.

\subsection{Empirical Measurement Procedure}
\label{sec:measurement}

Empirical measurements reported in
Table~\ref{tab:empirical_sizes_to_costs} were obtained by explicitly
materializing each fidelity in $\mathcal{F}$ and measuring its storage footprint
under realistic cloud storage configurations.
For each dataset, all visual representations were generated and uploaded to
object storage using standard cloud backends.

Amazon S3 was used as the primary reference platform due to its widespread
adoption and clearly defined tiered storage model (Standard, Standard--IA,
Glacier).
Object sizes were measured directly from stored objects via metadata queries
(e.g., \texttt{HEAD Object}), ensuring that compression and container overheads
were fully captured
\cite{aws_s3_storage_classes, aws_s3_pricing}.
We verified that relative size ratios and tier ordering are consistent with
Google Cloud Storage and Azure Blob Storage documentation
\cite{gcs_storage_classes, azure_blob_storage}.
Because \textsc{VOILA} depends only on relative acquisition costs, the resulting
normalized values are robust to provider-specific pricing and capture
deployment-invariant properties of tiered multimodal storage systems
\cite{hennessy2011computer}.

\begin{table}[t]
\centering
\small
\caption{\textbf{Empirical size measurements and derived costs (Edge--Cloud profile).}}
\label{tab:empirical_sizes_to_costs}
\begin{tabular}{lccccc}
\toprule
\textbf{Fidelity} &
\textbf{Avg. Size (KB)} &
$r(f)$ &
$b_{\mathrm{tier}}(f)$ &
$\tilde{c}(f)$ &
$c(f)$ \\
\midrule
Caption & 0.05 & $10^{-4}$ & 0.08 & 0.0800 & 7.6 \\
Resize $32{\times}32$ & 1.0 & 0.002 & 0.16 & 0.1601 & 15.1 \\
JPEG Q1 & 12 & 0.02 & 0.40 & 0.4012 & 37.8 \\
JPEG Q10 & 45 & 0.07 & 0.56 & 0.5642 & 53.2 \\
Full & 650 & 1.0 & 1.00 & 1.0600 & 100.0 \\
\bottomrule
\end{tabular}
\end{table}

\subsection{Sensitivity}
\label{sec:sensitivity}

VOILA depends only on relative cost differences.
Uniform rescaling or shifting of $c(f)$ leaves the selected fidelity unchanged
as long as higher fidelities remain strictly more expensive.
Empirically, VOILA’s accuracy--cost trade-offs remain stable across wide
variations in absolute cost magnitudes, confirming that the method captures
structural properties of multimodal systems rather than a specific pricing
scheme.
The key invariant is that \textsc{VOILA}’s routing behavior depends only on the \emph{ordering} of costs, not their absolute values: $c(\text{caption}) < c(\text{resize}) < c(\text{jpeg\_q1}) < c(\text{jpeg\_q10}) < c(\text{full})$ is preserved under any monotonic rescaling, and routing decisions are unchanged by uniform rescaling of all costs.
Deploying \textsc{VOILA} in a production system with measured per-tier latencies requires only recalibrating $\alpha$, $\beta$, $\gamma$ to match deployment-specific measurements; the framework supports this as a drop-in substitution.
